\textbf{Phần 1. Trắc nghiệm lựa chọn}
\setcounterpageref{ex}{0}

\begin{ex} %[3D3V33]
Nhiệt độ của cơ thể người bình thường và khỏe mạnh vào khoảng
\choice
{\True $37^\circ\text{C}$}
{$40^\circ\text{C}$}
{$35^\circ\text{C}$}
{$30^\circ\text{C}$}
\loigiai{
Nhiệt độ cơ thể người bình thường là khoảng $37^\circ\text{C}$.
\textbf{Đáp án của câu hỏi là A}
}
\end{ex}

\begin{ex} %[3D3V33]
Nhiệt độ nước đang sôi (ở điều kiện tiêu chuẩn) là
\choice
{$0^\circ\text{C}$}
{$50^\circ\text{C}$}
{\True $100^\circ\text{C}$}
{$80^\circ\text{C}$}
\loigiai{
Nước sôi ở nhiệt độ $100^\circ\text{C}$.
\textbf{Đáp án của câu hỏi là C}
}
\end{ex}

\begin{ex} %[3D3V33]
Buổi sáng nhiệt độ phòng là $24^\circ\text{C}$, đến buổi trưa nhiệt độ tăng thêm $6^\circ\text{C}$. Nhiệt độ buổi trưa là
\choice
{$28^\circ\text{C}$}
{\True $30^\circ\text{C}$}
{$32^\circ\text{C}$}
{$26^\circ\text{C}$}
\loigiai{
Nhiệt độ buổi trưa là: $24 + 6 = 30^\circ\text{C}$.
\textbf{Đáp án của câu hỏi là B}
}
\end{ex}

\begin{ex} %[3D3V33]
Dụng cụ dùng để đo nhiệt độ không khí hoặc nhiệt độ cơ thể người là
\choice
{Cân đồng hồ}
{Thước đo}
{Compa}
{\True Nhiệt kế}
\loigiai{
Dụng cụ chuẩn để đo nhiệt độ là nhiệt kế.
\textbf{Đáp án của câu hỏi là D}
}
\end{ex}

\begin{ex} %[3D3V33]
Nhiệt độ tại Sa Pa là $12^\circ\text{C}$, nhiệt độ tại Hà Nội là $26^\circ\text{C}$. Hỏi nhiệt độ tại Hà Nội cao hơn nhiệt độ tại Sa Pa bao nhiêu độ C?
\choice
{\True $14^\circ\text{C}$}
{$12^\circ\text{C}$}
{$16^\circ\text{C}$}
{$10^\circ\text{C}$}
\loigiai{
Chênh lệch nhiệt độ là: $26 - 12 = 14^\circ\text{C}$.
\textbf{Đáp án của câu hỏi là A}
}
\end{ex}

\begin{ex} %[3D3V33]
Nước đá (nước đóng băng) tan ở nhiệt độ nào?
\choice
{$10^\circ\text{C}$}
{$5^\circ\text{C}$}
{\True $0^\circ\text{C}$}
{$-10^\circ\text{C}$}
\loigiai{
Nước đá tan ở nhiệt độ $0^\circ\text{C}$.
\textbf{Đáp án của câu hỏi là C}
}
\end{ex}

\begin{ex} %[3D3V33]
Khi đo thân nhiệt cho bạn Nam, nhiệt kế chỉ $38^\circ\text{C} 5\text{ vạch nhỏ}$ (mỗi vạch nhỏ là $0,1^\circ\text{C}$, tức $38,5^\circ\text{C}$). Nhận xét nào sau đây đúng nhất?
\choice
{Bạn Nam hoàn toàn bình thường}
{\True Bạn Nam đang bị sốt nhẹ, cần nghỉ ngơi và theo dõi}
{Bạn Nam bị hạ thân nhiệt nguy hiểm}
{Bạn Nam đang bị đóng băng}
\loigiai{
Thân nhiệt trên $37,5^\circ\text{C}$ là có biểu hiện sốt.
\textbf{Đáp án của câu hỏi là B}
}
\end{ex}

\begin{ex} %[3D3V33]
Hôm qua nhiệt độ cao nhất là $32^\circ\text{C}$. Hôm nay do có không khí lạnh tràn về, nhiệt độ giảm đi $7^\circ\text{C}$. Nhiệt độ cao nhất hôm nay là
\choice
{$24^\circ\text{C}$}
{$26^\circ\text{C}$}
{$27^\circ\text{C}$}
{\True $25^\circ\text{C}$}
\loigiai{
Nhiệt độ hôm nay là: $32 - 7 = 25^\circ\text{C}$.
\textbf{Đáp án của câu hỏi là D}
}
\end{ex}

\begin{ex} %[3D3V33]
Nhiệt độ ngăn đá tủ lạnh thường duy trì ở mức
\choice
{\True Dưới $0^\circ\text{C}$ (âm độ C để làm đá)}
{Khoảng $15^\circ\text{C}$}
{Khoảng $25^\circ\text{C}$}
{Khoảng $50^\circ\text{C}$}
\loigiai{
Ngăn đá phải duy trì dưới $0^\circ\text{C}$ để nước đóng thành băng đá.
\textbf{Đáp án của câu hỏi là A}
}
\end{ex}

\begin{ex} %[3D3V33]
Nhiệt độ 4 ngày trong tuần lần lượt là: Thứ Hai: $25^\circ\text{C}$; Thứ Ba: $27^\circ\text{C}$; Thứ Tư: $23^\circ\text{C}$; Thứ Năm: $29^\circ\text{C}$. Ngày có nhiệt độ thấp nhất là
\choice
{Thứ Hai}
{Thứ Ba}
{\True Thứ Tư}
{Thứ Năm}
\loigiai{
Nhiệt độ thấp nhất là $23^\circ\text{C}$ rơi vào ngày Thứ Tư.
\textbf{Đáp án của câu hỏi là C}
}
\end{ex}

\begin{ex} %[3D3V33]
Trong một phòng thí nghiệm, dung dịch được đun từ $20^\circ\text{C}$ lên $80^\circ\text{C}$. Hỏi nhiệt độ dung dịch đã tăng bao nhiêu độ C?
\choice
{$50^\circ\text{C}$}
{\True $60^\circ\text{C}$}
{$70^\circ\text{C}$}
{$40^\circ\text{C}$}
\loigiai{
Nhiệt độ tăng thêm là: $80 - 20 = 60^\circ\text{C}$.
\textbf{Đáp án của câu hỏi là B}
}
\end{ex}

\begin{ex} %[3D3V33]
Nhiệt kế thủy ngân hoạt động dựa trên hiện tượng nào?
\choice
{Sự đổi màu của chất lỏng}
{Sự bốc hơi của nước}
{Trọng lực Trái Đất}
{\True Sự dãn nở vì nhiệt của chất lỏng}
\loigiai{
Khi nhiệt độ tăng, chất lỏng bên trong cột nhiệt kế nở ra và dâng cao lên.
\textbf{Đáp án của câu hỏi là D}
}
\end{ex}

\textbf{Phần 2. Trắc nghiệm đúng sai}
\setcounterpageref{ex}{0}

\begin{ex} %[3D3V33]
Xét tính đúng sai của các câu nói về nhiệt độ:
\choiceTF
{\True Đơn vị đo nhiệt độ thường dùng ở Việt Nam là độ C ($^\circ\text{C}$).}
{Nước đá tan ở $100^\circ\text{C}$.}
{\True Nhiệt độ cơ thể người khỏe mạnh bình thường khoảng $37^\circ\text{C}$.}
{Nhiệt độ ban đêm luôn cao hơn nhiệt độ ban ngày.}
\loigiai{
\begin{itemize}
    \item Ý a) Đúng: Kiến thức chuẩn.
    \item Ý b) Sai: Sai vì nước đá tan ở $0^\circ\text{C}$.
    \item Ý c) Đúng: Thực tế sinh học.
    \item Ý d) Sai: Sai vì ban đêm thường mát hơn ban ngày.
\end{itemize}
\begin{itemize}
    \item a) \textbf{Đúng}.
    \item b) \textbf{Sai}.
    \item c) \textbf{Đúng}.
    \item d) \textbf{Sai}.
\end{itemize}
}
\end{ex}

\begin{ex} %[3D3V33]
Dự báo thời tiết 3 thành phố: Hà Nội $28^\circ\text{C}$, Đà Nẵng $32^\circ\text{C}$, Đà Lạt $18^\circ\text{C}$:
\choiceTF
{\True Thành phố nóng nhất trong ba thành phố là Đà Nẵng.}
{\True Thành phố lạnh nhất trong ba thành phố là Đà Lạt.}
{Hà Nội lạnh hơn Đà Lạt $10^\circ\text{C}$.}
{\True Chênh lệch nhiệt độ giữa Đà Nẵng và Đà Lạt là $14^\circ\text{C}$.}
\loigiai{
\begin{itemize}
    \item Ý a) Đúng: $32^\circ\text{C}$ cao nhất.
    \item Ý b) Đúng: $18^\circ\text{C}$ thấp nhất.
    \item Ý c) Sai: Hà Nội nóng hơn Đà Lạt $28 - 18 = 10^\circ\text{C}$.
    \item Ý d) Đúng: $32 - 18 = 14^\circ\text{C}$.
\end{itemize}
\begin{itemize}
    \item a) \textbf{Đúng}.
    \item b) \textbf{Đúng}.
    \item c) \textbf{Sai}.
    \item d) \textbf{Đúng}.
\end{itemize}
}
\end{ex}

\textbf{Phần 3. Trả lời ngắn}
\setcounterpageref{ex}{0}

\begin{ex} %[3D3V33]
Buổi sáng nhiệt độ là $22^\circ\text{C}$, buổi trưa nhiệt độ tăng thêm $7^\circ\text{C}$. Tính nhiệt độ buổi trưa theo độ C.
\shortans{29}
\loigiai{
\textbf{Đáp số:} $29$\\
$22 + 7 = 29^\circ\text{C}$.
}
\end{ex}

\begin{ex} %[3D3V33]
Nhiệt độ nước đun trong ấm tăng từ $25^\circ\text{C}$ lên đến khi sôi ($100^\circ\text{C}$). Hỏi nhiệt độ đã tăng bao nhiêu độ C?
\shortans{75}
\loigiai{
\textbf{Đáp số:} $75$\\
$100 - 25 = 75^\circ\text{C}$.
}
\end{ex}

\begin{ex} %[3D3V33]
Thân nhiệt của em bé là $39^\circ\text{C}$. Sau khi uống thuốc hạ sốt, thân nhiệt giảm $2^\circ\text{C}$. Tính thân nhiệt của em bé sau khi hạ sốt.
\shortans{37}
\loigiai{
\textbf{Đáp số:} $37$\\
$39 - 2 = 37^\circ\text{C}$.
}
\end{ex}

\begin{ex} %[3D3V33]
Nhiệt độ phòng điều hòa là $24^\circ\text{C}$, nhiệt độ ngoài sân là $36^\circ\text{C}$. Hỏi ngoài sân nóng hơn trong phòng bao nhiêu độ C?
\shortans{12}
\loigiai{
\textbf{Đáp số:} $12$\\
$36 - 24 = 12^\circ\text{C}$.
}
\end{ex}

\begin{ex} %[3D3V33]
Nhiệt độ cao nhất trong ngày là $33^\circ\text{C}$, nhiệt độ thấp nhất trong ngày là $21^\circ\text{C}$. Tính mức chênh lệch nhiệt độ trong ngày theo độ C.
\shortans{12}
\loigiai{
\textbf{Đáp số:} $12$\\
$33 - 21 = 12^\circ\text{C}$.
}
\end{ex}

\begin{ex} %[3D3V33]
Nhiệt độ của nước trong chậu là $18^\circ\text{C}$. Người ta đổ thêm nước nóng làm nhiệt độ tăng gấp đôi. Hỏi nhiệt độ lúc này của nước trong chậu là bao nhiêu độ C?
\shortans{36}
\loigiai{
\textbf{Đáp số:} $36$\\
$18 \times 2 = 36^\circ\text{C}$.
}
\end{ex}
